\section{Вычислительный эксперимент}

Вычислительный эксперимент был направлен на иллюстрацию различий между тремя сценариями решения ECDLP, рассмотренными в работе: общим методом ро Полларда, методом Похлига--Хеллмана для групп с гладким порядком и атакой Смарта на аномальные кривые. Для каждого сценария подбирались такие входные данные, на которых соответствующий алгоритм можно было практически запустить и сопоставить с теоретическими выводами предыдущих разделов.

Для кривых без существенных особенностей (без уязвимостей ни к чему, кроме алгоритма ро Полларда) результаты сведены в таблицу~\ref{tab:pollard-rho}.

\begin{table}[htbp]
\centering
\begin{tabular}{|c|c|}
\hline
Размер кривой & Время вычисления \\
\hline
16 бит & $\sim 0{,}003$~с \\
32 бита & $\sim 0{,}5$~с \\
40 бит & $\sim 10$~с \\
64 бита & решение не найдено за 10 минут \\
\hline
\end{tabular}
\caption{Результаты для обычных кривых и метода ро Полларда}
\label{tab:pollard-rho}
\end{table}

Для кривых с гладким порядком (уязвимых к алгоритму Похлига--Хеллмана) результаты представлены в таблице~\ref{tab:pohlig-hellman}.

\begin{table}[htbp]
\centering
\begin{tabular}{|c|c|c|}
\hline
Размер кривой & Ограничение на множители & Время вычисления \\
\hline
64 бита & не более 16 бит & $\sim 0{,}005$~с \\
64 бита & не более 32 бит & $\sim 0{,}5$~с \\
128 бит & не более 32 бит & $\sim 0{,}6$~с \\
\hline
\end{tabular}
\caption{Результаты для кривых с гладким порядком и метода Похлига--Хеллмана}
\label{tab:pohlig-hellman}
\end{table}

Для аномальных кривых (к которым применима атака Смарта) соответствующие результаты приведены в таблице~\ref{tab:smart-attack}.

\begin{table}[htbp]
\centering
\begin{tabular}{|c|c|}
\hline
Размер кривой & Время вычисления \\
\hline
128 бит & $\sim 0{,}002$~с \\
256 бит & $\sim 0{,}01$~с \\
512 бит & $\sim 0{,}06$~с \\
\hline
\end{tabular}
\caption{Результаты для аномальных кривых и атаки Смарта}
\label{tab:smart-attack}
\end{table}

Следует также учитывать, что измерения для обычных кривых и для кривых с гладким порядком обладают высокой дисперсией. Это связано с вероятностной природой алгоритма ро Полларда, который как подзадача использовался и в алгоритме Похлига-Хеллмана. Поэтому приведённые времена следует понимать как характерные ориентировочные значения. В случае атаки Смарта разброс существенно меньше, поскольку здесь используется специальное структурное свойство аномальных кривых, а не общий вероятностный поиск коллизий.